% Target: CICM 2026 (19th Conference on Intelligent Computer Mathematics)
% Ljubljana, Slovenia, September 21-25, 2026
% Regular paper: up to 15 pages + bibliography, LNCS format
% Abstract deadline: March 25, 2026 | Paper deadline: April 1, 2026
% Published in: Springer LNAI
\documentclass[runningheads]{llncs}

\usepackage{amsmath,amssymb}
\usepackage{booktabs}
\usepackage{listings}
\usepackage{xcolor}
\usepackage{hyperref}
\usepackage{cleveref}
\usepackage{enumitem}

% Lean 4 listing style
\lstdefinelanguage{Lean4}{
  keywords={theorem, lemma, def, noncomputable, import, by, have, exact, rw,
            simp, ring, norm_num, apply, intro, induction, unfold, show,
            refine, suffices, convert, push_cast, field_simp, ext, linarith,
            private, where, with, match, if, then, else, let, in, do, return,
            structure, class, instance, open, section, namespace, end, variable,
            axiom, sorry, example, calc, fun},
  morekeywords=[2]{ℂ, ℝ, ℕ, ℤ},
  sensitive=true,
  morecomment=[l]{--},
  morecomment=[s]{/-}{-/},
  morestring=[b]",
}

\lstset{
  language=Lean4,
  basicstyle=\ttfamily\small,
  keywordstyle=\color{blue}\bfseries,
  keywordstyle=[2]\color{purple},
  commentstyle=\color{gray}\itshape,
  stringstyle=\color{red},
  breaklines=true,
  columns=flexible,
  frame=single,
  xleftmargin=2em,
  numbers=left,
  numberstyle=\tiny\color{gray},
  numbersep=5pt,
  literate=
    {ℝ}{{$\mathbb{R}$}}1
    {ℂ}{{$\mathbb{C}$}}1
    {ℕ}{{$\mathbb{N}$}}1
    {ℤ}{{$\mathbb{Z}$}}1
    {↑}{{$\uparrow$}}1
    {→}{{$\rightarrow$}}1
    {←}{{$\leftarrow$}}1
    {⟨}{{$\langle$}}1
    {⟩}{{$\rangle$}}1
    {∀}{{$\forall$}}1
    {∃}{{$\exists$}}1
    {λ}{{$\lambda$}}1,
}

\begin{document}

\title{Formal Verification of Alternative Mathematical Frameworks:\\
A Case Study Using Lean~4}
\titlerunning{Formal Verification of Alternative Mathematical Frameworks}

\author{Ian M. Arnoldy}
\institute{Independent Researcher}

\maketitle

\begin{abstract}
We present a systematic protocol for applying interactive theorem provers to
claims originating from alternative or non-standard mathematical frameworks.
Using Lean~4~\cite{demoura2021lean4} and mathlib~\cite{mathlib}, we extract formalizable assertions from such sources,
translate them into dependent type theory, and attempt proof or disproof---reporting
results without prejudice. We demonstrate the protocol through a case study on
Eric Dollard's ``versor algebra,'' a non-standard extension of electrical
engineering mathematics. Results are deliberately mixed: we verify a $\mathbb{Z}_4$
cyclic group structure (trivially correct), disprove a claimed versor-form
equivalence (sign error), establish by algebraic necessity that the proposed
operators collapse to standard 4th roots of unity, and trace the algebraic path
from $\mathbb{Z}_4$ through Clifford algebras $\mathrm{Cl}(1,1)$,
$\mathrm{Cl}(3,0)$, $\mathrm{Cl}(1,3)$ to spacetime algebra. We propose a
five-category taxonomy for classifying claims from alternative frameworks:
\textsc{Verified}, \textsc{Disproved}, \textsc{Algebraic~Necessity},
\textsc{Ambiguous}, and \textsc{Unfalsifiable}. To our knowledge, this represents
the first application of interactive theorem provers to formalize claims from
non-mainstream mathematical sources.
\end{abstract}

\keywords{formal verification \and Lean~4 \and alternative mathematics \and
Clifford algebras \and versor algebra}

% ============================================================================
\section{Introduction}
\label{sec:intro}
% ============================================================================

Interactive theorem provers have verified some of mathematics' most challenging
results: Hales' Kepler conjecture via the Flyspeck project~\cite{hales2017formal},
the Four Color Theorem~\cite{gonthier2008four}, Scholze's liquid tensor
experiment~\cite{commelin2022liquid}, the polynomial Freiman-Ruzsa
conjecture~\cite{tao2023pfr}, and the consistency of Quine's New
Foundations~\cite{wilshaw2024nf}. In each case, the target was mainstream
mathematics---contested, complex, or uncertain, but arising within accepted
mathematical frameworks.

A conspicuous gap remains: no prior work applies interactive theorem provers to
claims from \emph{alternative} or \emph{non-standard} mathematical frameworks.
By ``alternative framework'' we mean a coherent system of mathematical claims
originating outside peer-reviewed mathematics, typically from independent
researchers, that proposes non-standard definitions, operations, or equivalences.
Such frameworks are distinct from non-standard \emph{foundations} (e.g., HoTT,
constructive mathematics), which are rigorous alternative approaches proposed
within the mathematical community~\cite{shulman2024strange}.

This gap is structural, not accidental. Four factors conspire to maintain it:
\begin{enumerate}
\item \textbf{Cultural separation.} Formal verification and alternative
  mathematics communities have zero overlap.
\item \textbf{No academic incentive.} Formalizing alternative claims yields no
  new theorems and no library contributions, while risking the appearance of
  legitimizing pseudoscience.
\item \textbf{Asymmetric difficulty.} Translating informal non-standard claims
  into dependent type theory requires both mathematical expertise and judgment
  about ambiguous notation.
\item \textbf{No client.} Proponents of alternative frameworks rarely seek
  formal verification (which might disprove claims), and mainstream mathematicians
  see no need (they already reject such claims on conventional grounds).
\end{enumerate}

We argue that this gap represents a missed opportunity. Formal verification of
alternative frameworks is methodologically interesting precisely when it produces
\emph{mixed} results: some claims verified, others disproved, still others
classified as ambiguous or unfalsifiable. Pure debunking is uninteresting (it
offers no methodological contribution), and pure verification is unlikely for
frameworks that diverge from standard mathematics. Mixed results demonstrate the
discriminating power of formal methods and reveal genuine mathematical structure
worth analyzing.

\paragraph{Contributions.} This paper makes four contributions:
\begin{enumerate}
\item A six-step \emph{protocol} for extracting, translating, and verifying
  claims from alternative mathematical frameworks (\Cref{sec:protocol}).
\item A \emph{case study} applying the protocol to Eric Dollard's versor
  algebra, producing over 200 formally verified theorems across 8 Lean~4 proof files
  with zero \texttt{sorry} gaps (\Cref{sec:case-study}); the full repository
  extends to 59 proof files with over 2,537 verified declarations, including
  four Lie algebras with certified mathlib \texttt{LieRing}/\texttt{LieAlgebra}
  instances.
\item A five-category \emph{taxonomy} for classifying claims from alternative
  frameworks: \textsc{Verified}, \textsc{Disproved},
  \textsc{Algebraic~Necessity}, \textsc{Ambiguous}, and \textsc{Unfalsifiable}
  (\Cref{sec:taxonomy}).
\item A constructive \emph{algebraic hierarchy} tracing the path from the
  source claims through Clifford algebras to spacetime algebra, demonstrating
  that formal methods can suggest corrections, not merely judge
  (\Cref{sec:hierarchy}).
\end{enumerate}

% ============================================================================
\section{Related Work}
\label{sec:related}
% ============================================================================

\subsection{Formalization of Contested Mathematics}

The closest parallels to our work involve formalizations where the mathematical
community itself had doubts about correctness. The Flyspeck
project~\cite{hales2017formal} verified Hales' proof of the Kepler conjecture
after referees could not fully evaluate it. The liquid tensor
experiment~\cite{commelin2022liquid} verified a theorem where even the author
(Scholze) had reservations. Wilshaw's formalization of NF
consistency~\cite{wilshaw2024nf} addressed a proof that mainstream reviewers
could not confidently evaluate. In each case, the target was
contested-within-mainstream, not alternative/fringe.

\subsection{Formalization of Physical Theories}

Tooby-Smith's PhysLean/HepLean project~\cite{toobysmith2025physlean}
digitalizes mainstream physics in Lean~4, from classical mechanics through
quantum field theory. Bobbin et al.~\cite{bobbin2024chemical} formalized
Langmuir and BET adsorption theories, demonstrating that Lean catches faulty
logic and reveals hidden assumptions---the closest methodological parallel to
our work, though applied to mainstream physical chemistry.

\subsection{Historical Mathematics Formalization}

Fleuriot~\cite{fleuriot2001newton} formalized Newton's geometric
proofs in Isabelle/HOL and discovered a flaw in one of Newton's proofs. This
is the closest near-miss to our work: a proof assistant applied to historical
non-rigorous reasoning that finds errors. The distinction is that Newton is
mainstream; the ``non-standard'' aspect is proof style, not framework.

\subsection{Clifford Algebra Formalization}

Wieser and Song~\cite{wieser2021clifford} formalized Clifford algebras in
Lean~3/mathlib, providing the mathematical objects that overlap with our case
study. Their work formalizes standard theory; ours applies it to non-standard
claims.

\subsection{Philosophy of Formal Verification}

Avigad~\cite{avigad2021understanding} examines how formal verification
contributes to mathematical understanding. Imbert and
Ardourel~\cite{imbert2023formal} analyze formal verification's epistemological
role in computational science. Shulman~\cite{shulman2024strange} discusses proof
assistants for exploring alternative foundations (HoTT), though these are
rigorous alternatives proposed by credentialed mathematicians.

% ============================================================================
\section{The Protocol}
\label{sec:protocol}
% ============================================================================

We propose a six-step protocol for applying interactive theorem provers to
alternative mathematical frameworks. The individual techniques at each step are
standard; the contribution is the systematic application to a novel domain.

\begin{enumerate}
\item \textbf{Source Identification.} Identify an alternative mathematical
  framework containing specific algebraic or structural claims amenable to
  formalization. The framework must make assertions beyond rhetorical or
  philosophical claims.

\item \textbf{Claim Extraction.} Enumerate all formalizable assertions from the
  source material. Separate mathematical claims (``$j^4 = 1$'') from rhetoric
  (``electricity is fundamental''), physics claims (``aether exists''), and
  philosophical framing. Document the extraction process, including claims
  deliberately excluded and why.

\item \textbf{Translation.} Map each extracted claim to dependent type theory.
  Document all translation choices, especially where source notation is ambiguous.
  When multiple interpretations exist, formalize the most charitable one
  (\emph{principle of charity}).

\item \textbf{Verification Attempt.} For each translated claim, attempt proof
  or disproof using standard tactics. Record the outcome: proved, disproved
  (with counterexample or contradiction), or neither (claim may be ambiguous
  or unformalizable).

\item \textbf{Taxonomy Classification.} Assign each claim to one of five
  categories (\Cref{sec:taxonomy}).

\item \textbf{Neutral Reporting.} Present all results without editorial judgment
  on the source framework. The role of the formal verifier is analogous to that
  of an auditor: to report findings, not to advocate or condemn.
\end{enumerate}

% ============================================================================
\section{Case Study: Dollard's Versor Algebra}
\label{sec:case-study}
% ============================================================================

\subsection{Background}

Eric Dollard is an independent researcher in electrical engineering who has
proposed a ``versor algebra'' as an extension of standard complex-number methods
for analyzing polyphase power systems~\cite{fortescue1918} and transmission
lines~\cite{dollard-versor,dollard-fourquad}. The framework introduces three
operators: $j$ (the standard imaginary unit), $h$ (a proposed ``hyperbolic''
operator with $h^2 = 1$, $h^1 = -1$), and $k = hj$ (a ``sequence'' operator).
Dollard claims these operators form a novel algebraic structure with applications
to electromagnetic theory.

The source materials are four self-published documents: \emph{Lone Pine
Writings}, \emph{Four Quadrant Theory}, \emph{Versor Algebra}, and
\emph{Advanced Versor Algebra II-23}~\cite{dollard-versor,dollard-fourquad}.
These documents are not peer-reviewed and are not widely available; we
therefore reproduce the core definitions here for self-containment.
We extracted 17 formalizable claims from these sources.

\paragraph{Dollard's axioms.}
Dollard postulates three operators $j$, $h$, $k$ in a unital associative
algebra, subject to five stated properties:
\begin{enumerate}[label=(A\arabic*),nosep]
  \item $j^2 = -1$ \quad (the standard imaginary unit),
  \item $hj = k$ \quad (product definition of the sequence operator),
  \item $jk = 1$ \quad (cancellation),
  \item $h^1 = -1$ \quad ($h$ evaluates to $-1$),
  \item $h^2 = 1$ \quad ($h$ squared is unity).
\end{enumerate}
Dollard writes $h = \sqrt{+1}^{1/2}$, using both the square root and a
fractional exponent, which is mathematically ambiguous
(see \Cref{sec:algnec}).  The intended interpretation appears to be an
operator that squares to $+1$ but is itself $-1$---distinct from the
principal square root.

\paragraph{Sign convention.} A critical detail: Dollard uses the admittance
convention $Y = G - jB$ (negative susceptance), not the standard $Y = G + jB$.
This affects the telegraph equation expansion. Our formalization verifies both
conventions for comparison.

\subsection{\textsc{Verified}: $\mathbb{Z}_4$ Cyclic Group Structure}

Dollard's operators $\{1, j, h, k\}$ with $j^2 = -1$, $h = -1$, $k = hj = -j$,
and $jk = 1$ form a group under multiplication. This group is isomorphic to
$\mathbb{Z}_4$, the cyclic group of order 4, which is also the group of 4th
roots of unity $\{1, i, -1, -i\}$ under complex multiplication.

\begin{lstlisting}[caption={Versor operators as 4th roots of unity (excerpt from \texttt{basic\_operators.lean})}]
def j : ℂ := Complex.I
def h : ℂ := -1
def k : ℂ := h * j

theorem j_fourth_power : j ^ 4 = 1 := by ...
theorem h_squared : h ^ 2 = 1 := by ...
theorem jk_cancellation : j * k = 1 := by ...
\end{lstlisting}

All 16 products of the Cayley table are verified. The structure is trivially
correct---it is standard complex arithmetic. The verification confirms Dollard's
framework includes correct standard results, though the algebraic content is
not novel.

\subsection{\textsc{Disproved}: Versor Form Equivalence}
\label{sec:disproof}

Dollard claims that the telegraph equation product $ZY = (R + jX)(G - jB)$ can
be equivalently written in ``versor form'' as:
\[
  ZY = h(XB + RG) + j(XG - RB)
\]

With $h = -1$ (as established by the $\mathbb{Z}_4$ structure):
\begin{align*}
  \text{Versor:} \quad & h(XB{+}RG) + j(XG{-}RB) = -(XB{+}RG) + j(XG{-}RB) \\
  \text{Telegraph:} \quad & (RG{+}XB) + j(XG{-}RB)
\end{align*}

The imaginary parts match, but the real parts differ by sign. The forms are equal
only when $RG + XB = 0$---a degenerate case where both dissipation and storage
vanish simultaneously.

\begin{lstlisting}[caption={Versor form disproof (excerpt from \texttt{telegraph\_equation.lean})},basicstyle=\ttfamily\footnotesize]
theorem versor_form_value (R X G B : ℝ) :
    h * (↑(factor_XB X B + factor_RG R G) : ℂ) +
    j * (↑(factor_XG X G - factor_RB R B) : ℂ) =
    Complex.mk (-(factor_XB X B + factor_RG R G))
               (factor_XG X G - factor_RB R B) := by
  simp only [h, j]
  apply Complex.ext <;> simp [...]
\end{lstlisting}

This disproof is robust: it holds regardless of sign convention (Dollard or
standard), because $h = -1$ negates the real part in either case. Furthermore,
the versor form is \emph{irreparable}: we prove that the coefficient on the
real part must equal 1 for equivalence (\texttt{versor\_repaired\_with\_one}),
but $h \neq 1$ (\texttt{h\_ne\_one}).

\subsection{\textsc{Algebraic Necessity}: $h = -1$ Forced}
\label{sec:algnec}

The most structurally interesting result: Dollard's three core axioms
$j^2 = -1$, $hj = k$, $jk = 1$ \emph{alone} force $h = -1$ over any field.
The proof is elementary:
\[
  jk = 1 \implies k = j^{-1} = -j, \qquad hj = k = -j \implies h = -1
\]

The axioms $h^1 = -1$ and $h^2 = 1$ are redundant---they follow from the three
core axioms. This means Dollard's algebra has zero degrees of freedom: any
unital associative algebra satisfying the core axioms is forced to
$\mathbb{Z}_4$.

\begin{lstlisting}[caption={Algebraic necessity over any field (excerpt from \texttt{algebraic\_necessity.lean})}]
-- Section variables: {F : Type*} [Field F] {j' h' k' : F}
theorem algebraic_necessity_general
    (hj2 : j' ^ 2 = -1) (hhj : h' * j' = k')
    (hjk : j' * k' = 1) :
    h' = -1 := by ...
\end{lstlisting}

This result closes the ``rehabilitation question'': to obtain a non-trivial $h$
(e.g., in $\mathrm{Cl}(1,1)$ where $e_1^2 = 1$ but $e_1 \neq \pm 1$), one
must drop at least two of Dollard's three core axioms. The resulting algebra
is a \emph{different} system, not a repair of Dollard's.

\paragraph{The $D_8$ alternative.}
Wilson~\cite{wilson-private} observes that versors in the standard
Clifford-algebraic sense satisfy anticommutativity: in $\mathrm{Cl}(1,1)$,
generators $e_1, e_2$ with $e_1^2 = +1$, $e_2^2 = -1$ satisfy
$e_1 e_2 = -e_2 e_1$. Identifying $h = e_1$ and $j = e_2$, the group
generated by $\{j, h\}$ has 8 elements $\{\pm 1, \pm j, \pm h, \pm jh\}$
and is isomorphic to the dihedral group $D_8$.
Our algebraic necessity result confirms this from the complementary direction:
Dollard's (implicit) commutativity assumption forces $h = -1$, collapsing
$D_8$ to $\mathbb{Z}_4$. The sign error identified in \Cref{sec:disproof}
is a direct consequence of this collapse---had Dollard worked in
$\mathrm{Cl}(1,1)$ with anticommuting generators, the versor form would
require different algebraic treatment entirely.

\subsection{\textsc{Ambiguous}: Notation Issues}

Dollard writes $h = \sqrt{+1}^{1/2}$, which is mathematically ambiguous.
In $\mathbb{R}$, $\sqrt{1} = 1$ and $1^{1/2} = 1$, so $h = 1$---but Dollard
claims $h = -1$. The notation admits multiple interpretations, none of which
are standard.

Formal methods precisely \emph{identify} this ambiguity as a contribution:
rather than resolving it by fiat, we show that the algebraic necessity result
(\Cref{sec:algnec}) makes the notation irrelevant. Whatever
$\sqrt{+1}^{1/2}$ means, $h$ must equal $-1$ in any algebra satisfying
Dollard's core axioms.

\subsection{\textsc{Unfalsifiable}: Physics Claims}

Dollard defines $\Psi$ (total dielectric induction, in Coulombs) and $\Phi$
(total magnetic induction, in Webers), with $Q = \Psi \cdot \Phi$ and
$W = dQ/dt$. Dimensionally, $\text{Coulomb} \times \text{Weber} =
\text{J} \cdot \text{s}$ (action), making $Q$ the action and $W$ the energy.
This is confirmed by primary source analysis: Dollard's $d\Psi/dt = I$
(current) and $d\Phi/dt = E$ (voltage) exactly reproduce the Lagrangian
circuit variables where charge $q$ is the generalized coordinate and flux
$\lambda$ is the conjugate momentum. The dimensional
analysis is verified in Lean~4 (\texttt{circuit\_action.lean}, 0~\texttt{sorry}).
The framework is dimensionally correct but not novel---it is Lagrangian circuit
theory, known since Cherry~\cite{cherry1951}.

Several remaining claims---``electricity is fundamental,''
``aether exists as a fifth state of matter,'' ``$W = dQ/dt$ replaces
$E = mc^2$''---are not addressable by formal mathematical methods. We classify
these as \textsc{Unfalsifiable} (by theorem provers), noting that this
classification is itself a useful result: it separates the checkable from the
uncheckable.

% ============================================================================
\section{A Taxonomy for Claims from Alternative Frameworks}
\label{sec:taxonomy}
% ============================================================================

Based on the case study, we propose a five-category taxonomy for classifying
claims from alternative mathematical frameworks:

\begin{table}[ht]
\centering
\caption{Five-category taxonomy for claims from alternative frameworks}
\label{tab:taxonomy}
\begin{tabular}{@{}lp{4.5cm}p{4cm}@{}}
\toprule
\textbf{Category} & \textbf{Definition} & \textbf{Example} \\
\midrule
\textsc{Verified} & Claim is mathematically correct; formal proof produced &
  $\mathbb{Z}_4$ group structure \\
\textsc{Disproved} & Claim is mathematically false; formal counterexample or
  contradiction produced & Versor form equivalence \\
\textsc{Algebraic Necessity} & Premises force a conclusion the author did not
  intend & $h = -1$ forced by core axioms \\
\textsc{Ambiguous} & Claim cannot be formalized without resolving
  notation/definition gaps & $h = \sqrt{+1}^{1/2}$ \\
\textsc{Unfalsifiable} & Claim is not addressable by formal methods &
  Physics/empirical claims \\
\bottomrule
\end{tabular}
\end{table}

We argue that mixed results across these categories are the \emph{expected} and
methodologically interesting case. Uniform rejection would suggest the framework
is too incoherent to formalize. Uniform verification would suggest the framework
is merely standard mathematics with non-standard notation. Mixed results---some
verified, some disproved, some ambiguous---reveal genuine structure worth
analyzing and demonstrate the discriminating power of formal methods.

The \textsc{Algebraic Necessity} category merits special attention. It captures
cases where the author's axioms entail consequences they did not recognize or
intend. This phenomenon has precedents across several fields: as a consequence
of Hilbert's implicit definitions (the Frege-Hilbert controversy), as
``implicit commitment'' in epistemic logic~\cite{fagin1988belief}, and as
overspecification in formal methods. The Banach-Tarski paradox is a canonical
example at larger scale---the axiom of choice forces consequences its adopters
did not foresee.

We note a duality with Lakatos's \emph{Proofs and
Refutations}~\cite{lakatos1976}: Lakatos addresses axiom systems that are too
\emph{weak}, admitting unintended models (``monsters'' requiring barring). Our
\textsc{Algebraic Necessity} addresses axiom systems that are too \emph{strong},
forcing unintended consequences. These are dual failure modes:

\begin{center}
\small
\begin{tabular}{@{}llll@{}}
\toprule
\textbf{Failure} & \textbf{Direction} & \textbf{Lakatos} & \textbf{This work} \\
\midrule
Axioms too weak & Admit unintended models & Monster-barring & --- \\
Axioms too strong & Force unintended conclusions & --- & \textsc{Alg.\ Necessity} \\
\bottomrule
\end{tabular}
\end{center}

This duality suggests that \textsc{Algebraic Necessity} fills a genuine gap in
the vocabulary for analyzing axiom systems---particularly relevant when formal
methods are applied to frameworks whose authors may not have fully explored the
consequences of their axioms.

% ============================================================================
\section{The Clifford Algebra Hierarchy}
\label{sec:hierarchy}
% ============================================================================

Beyond classification, formal methods can suggest \emph{corrections}. We trace
the algebraic path from Dollard's $\mathbb{Z}_4$ to the spacetime algebra where
electromagnetism becomes geometry:

\begin{table}[ht]
\centering
\caption{Algebraic hierarchy from $\mathbb{Z}_4$ to spacetime algebra}
\label{tab:hierarchy}
\begin{tabular}{@{}clcl@{}}
\toprule
\textbf{Level} & \textbf{Algebra} & \textbf{Dim} & \textbf{Physical Content} \\
\midrule
1 & $\mathbb{Z}_4$ & 4 & Versor operators (trivial, forced by axioms) \\
2 & $\mathrm{Cl}(1,1)$ & 4 & Idempotent wave decomposition \\
3 & $\mathrm{Cl}(3,0)$ & 8 & Pauli algebra; $F = E + IB$ \\
4 & $\mathrm{Cl}(1,3)$ & 16 & Spacetime algebra; $\nabla F = J$ \\
\bottomrule
\end{tabular}
\end{table}

Each level is formalized in Lean~4 with zero \texttt{sorry} gaps.

At Level~2, $\mathrm{Cl}(1,1)$ provides what Dollard was reaching for: a
non-trivial operator $e_1$ with $e_1^2 = +1$ but $e_1 \neq \pm 1$. However,
this algebra is non-commutative ($e_1 e_2 \neq e_2 e_1$), and Dollard's axiom
$jk = 1$ fails ($e_2 \cdot e_{12} = e_1 \neq 1$). The algebra offers idempotent
projectors $P_{\pm} = \frac{1 \pm e_1}{2}$ that decompose signals into
forward and backward wave components---a genuine mathematical tool for
transmission line analysis.

At Level~4, $\mathrm{Cl}(1,3)$ is the spacetime algebra~\cite{hestenes1966,doran2003}
where Maxwell's four equations collapse to one: $\nabla F = J$. This is the genuine unified field
framework for classical electromagnetism, achieved through geometric algebra
rather than through Dollard's commutative $\mathbb{Z}_4$.

The hierarchy demonstrates that Dollard's \emph{instinct}---algebraic
unification of electromagnetic phenomena---was correct, while his
\emph{execution}---a commutative group with four elements---was far too
restrictive. Formal methods can articulate this distinction precisely.

% ============================================================================
\section{Discussion}
\label{sec:discussion}
% ============================================================================

\subsection{Why Mixed Results Matter}

The case study produced all five taxonomy categories across 17 claims. This
distribution validates the protocol's discriminating power: it is neither
a rubber stamp nor a blanket dismissal.

The sign error in the versor form (\Cref{sec:disproof}) could not have been
established with certainty without formalization. Informal mathematical
reasoning might notice the discrepancy, but Lean provides machine-checked
certainty---and proves the error is irreparable.

\subsection{Why the Gap Exists}

The four structural factors identified in \Cref{sec:intro}---cultural
separation, no academic incentive, asymmetric difficulty, no client---explain
why this application domain has been overlooked. We note that none of these
factors are \emph{mathematical}; the gap is sociological. The mathematical
tools are entirely adequate: standard Lean~4 tactics suffice for all proofs in
our case study, and no custom tactics or library extensions were required.

\subsection{The Constructive Dimension}

Our Clifford algebra hierarchy (\Cref{sec:hierarchy}) goes beyond demarcation
to \emph{mathematical rehabilitation}: it traces the path from the alternative
claim to the correct mathematical framework. This constructive contribution---
showing not just what is wrong but what the author \emph{should} have
said---may be the most valuable aspect of applying formal methods to
alternative frameworks.

\subsection{Limitations}

Our protocol is demonstrated on a single case study; generalizability to other
alternative frameworks is conjectured but not proven. Dollard's source materials
are self-published presentations, not peer-reviewed papers, introducing
translation ambiguities that are inherent to the domain. All Lean~4 proofs
compile against mathlib (the full 59-file repository builds with
zero errors, zero \texttt{sorry} gaps). We formalize algebraic and structural
claims only; we make no claims about the physics, engineering utility, or
broader validity of Dollard's work.

% ============================================================================
\section{Conclusion}
\label{sec:conclusion}
% ============================================================================

We have demonstrated that interactive theorem provers can be productively
applied to claims from alternative mathematical frameworks. The protocol
produces mixed results---verified, disproved, algebraically necessary,
ambiguous, and unfalsifiable---that reveal genuine mathematical structure.
The five-category taxonomy provides a systematic classification scheme.
The Clifford algebra hierarchy shows that formal methods can constructively
suggest the ``correct'' version of an alternative framework's aspirations.

The value of this approach lies not in judgment but in discrimination: formal
methods separate what is mathematically checkable from what is not, what is
correct from what is wrong, and what is novel from what is standard. We call
for further case studies on other alternative mathematical frameworks to
validate and refine the protocol.

The case study presented here covers the 8 proof files most directly tied to
the Dollard verification. The full repository extends the algebraic hierarchy
through grand unified theories (SO(10), SO(14)) and exceptional algebras
(E$_8$), totaling 59 proof files with over 2,537 verified declarations---all
compiling with zero \texttt{sorry} gaps. Recent work includes mathlib
\texttt{LieRing} and \texttt{LieAlgebra~$\mathbb{R}$} typeclass instances for
four Lie algebras (so(1,3), sl(3), sl(5), so(10)), connecting the hand-built
structures to mathlib's abstract type system, and a machine-verified resolution
of the three-generation problem via $E_8 \supset \mathrm{SU}(9)/\mathbb{Z}_3$.
Source code, source materials, and analysis documents are available from
the corresponding author upon reasonable request.

% ============================================================================
\section*{Appendix: Complete Claim Summary}
\label{sec:appendix}
% ============================================================================

\begin{table}[ht]
\centering
\caption{All 17 formalized claims with taxonomy classification}
\label{tab:claims}
\small
\begin{tabular}{@{}clll@{}}
\toprule
\textbf{\#} & \textbf{Claim} & \textbf{Category} & \textbf{Proof File} \\
\midrule
1 & $\{1,j,h,k\}$ as 4th roots of unity & \textsc{Verified} & \texttt{basic\_operators} \\
2 & Telegraph expansion $(R+jX)(G-jB)$ & \textsc{Verified} & \texttt{telegraph\_equation} \\
3 & Four-factor independence & \textsc{Disproved}$^*$ & \texttt{telegraph\_equation} \\
4 & Versor form equivalence & \textsc{Disproved} & \texttt{telegraph\_equation} \\
5 & Universal polyphase formula & \textsc{Verified} & \texttt{polyphase\_formula} \\
6 & $h = \sqrt{+1}^{1/2}$ notation & \textsc{Ambiguous} & \texttt{algebraic\_necessity} \\
7 & Sequence periodicity & \textsc{Verified} & \texttt{basic\_operators} \\
8 & $Q = \Psi \cdot \Phi$ (action) & \textsc{Verified}$^\dagger$ & \texttt{circuit\_action} \\
9 & Electricity as fundamental & \textsc{Unfalsifiable} & --- \\
10 & Aether as 5th state & \textsc{Unfalsifiable} & --- \\
11 & $W = dQ/dt$ replaces $E = mc^2$ & \textsc{Unfalsifiable} & --- \\
12 & Over-unity energy & \textsc{Unfalsifiable} & --- \\
13 & Dielectric counterspace & \textsc{Unfalsifiable} & --- \\
14 & Quaternary expansion & \textsc{Verified}$^\ddagger$ & (computational) \\
15 & Fortescue is domain-specific & \textsc{Verified} & (empirical) \\
16 & $A_4$ matches versor matrix & \textsc{Verified} & (computational) \\
17 & $jk=1$ forces $h=-1$ & \textsc{Alg.\ Necessity} & \texttt{algebraic\_necessity} \\
\bottomrule
\end{tabular}

\vspace{0.5em}
{\footnotesize
$^*$Four terms exist but satisfy $(RG)(XB) = (XG)(RB)$; not independent. \\
$^\dagger$Dimensionally correct (Weber$\cdot$Coulomb = J$\cdot$s) but not novel (Cherry, 1951). \\
$^\ddagger$Mathematically correct but practically useless (deprecated).
}
\end{table}

\paragraph{Note on mathlib.} Mathlib includes a comprehensive formalization of
Clifford algebras (\texttt{Mathlib.LinearAlgebra.CliffordAlgebra})~\cite{wieser2021clifford}
with grading, spin groups, and conjugation. Our case study uses explicit
finite-dimensional representations for pedagogical clarity, but a natural
extension would unify our ad-hoc $\mathrm{Cl}(p,q)$ types with mathlib's
abstract \texttt{CliffordAlgebra} via isomorphism proofs. As a step in this
direction, four Lie algebras in the extended repository---so(1,3), sl(3),
sl(5), and so(10)---now carry certified \texttt{LieRing} and
\texttt{LieAlgebra~$\mathbb{R}$} instances via
\texttt{Mathlib.Algebra.Lie.Basic}, connecting the explicit constructions to
mathlib's abstract type system.

% ============================================================================
\section*{Acknowledgements}
% ============================================================================

The machine verification was performed using Lean~4 and the mathlib
library. The author thanks the Lean and mathlib communities for providing
the infrastructure that made this work possible.
R.~A.~Wilson provided valuable feedback on the relationship between
Dollard's commutative versor algebra and the dihedral group $D_8$
arising from the standard anticommutative Clifford structure.
Proof engineering and manuscript preparation were assisted by Claude
(Anthropic); all mathematical content, research direction, and
interpretation are the sole responsibility of the author.

% ============================================================================
% References
% ============================================================================

\begin{thebibliography}{23}

\bibitem{hales2017formal}
Hales, T., et al.: A formal proof of the Kepler conjecture.
Forum of Mathematics, Pi \textbf{5}, e2 (2017)

\bibitem{gonthier2008four}
Gonthier, G.: Formal proof---The Four-Color Theorem.
Notices of the AMS \textbf{55}(11), 1382--1393 (2008)

\bibitem{commelin2022liquid}
Commelin, J., et al.: Completion of the Liquid Tensor Experiment.
Lean Community (2022)

\bibitem{tao2023pfr}
Tao, T., et al.: Polynomial Freiman-Ruzsa conjecture.
Lean~4 formalization (2023)

\bibitem{wilshaw2024nf}
Wilshaw, S.: New Foundations is consistent.
Lean formalization, leanprover-community/con-nf (2024)

\bibitem{shulman2024strange}
Shulman, M.: Strange new universes: Proof assistants and synthetic foundations.
Bull.\ AMS \textbf{61}(2) (2024)

\bibitem{toobysmith2025physlean}
Tooby-Smith, J.: HepLean: Digitalising high energy physics.
Computer Physics Communications \textbf{308} (2025)

\bibitem{bobbin2024chemical}
Bobbin, M.P., et al.: Formalizing chemical physics using the Lean theorem prover.
Digital Discovery, RSC (2024)

\bibitem{fleuriot2001newton}
Fleuriot, J.: A Combination of Geometry Theorem Proving and Nonstandard
Analysis with Application to Newton's Principia. Springer (2001)

\bibitem{wieser2021clifford}
Wieser, E., Song, U.: Formalizing Geometric Algebra in Lean.
Advances in Applied Clifford Algebras \textbf{31} (2021)

\bibitem{avigad2021understanding}
Avigad, J.: Understanding, formal verification, and the philosophy of mathematics.
Journal of the Indian Council of Philosophical Research (2021)

\bibitem{imbert2023formal}
Imbert, C., Ardourel, V.: Formal verification, scientific code, and
epistemological heterogeneity. Philosophy of Science (2023)

\bibitem{demoura2021lean4}
de Moura, L., Ullrich, S.: The Lean~4 theorem prover and programming language.
CADE-28, LNCS \textbf{12699} (2021)

\bibitem{mathlib}
The mathlib community: The Lean mathematical library. CPP 2020.

\bibitem{fortescue1918}
Fortescue, C.L.: Method of symmetrical co-ordinates applied to the solution
of polyphase networks. AIEE Transactions \textbf{37}(2), 1027--1140 (1918)

\bibitem{hestenes1966}
Hestenes, D.: Space-Time Algebra. Gordon and Breach (1966)

\bibitem{doran2003}
Doran, C., Lasenby, A.: Geometric Algebra for Physicists.
Cambridge University Press (2003)

\bibitem{dollard-versor}
Dollard, E.P.: Versor Algebra: As Applied to Polyphase Power Systems.
Self-published presentation.

\bibitem{dollard-fourquad}
Dollard, E.P.: Four Quadrant Theory: Advanced Electrical Applications.
Self-published presentation.

\bibitem{cherry1951}
Cherry, E.C.: Some general theorems for non-linear systems possessing reactance.
The London, Edinburgh, and Dublin Philosophical Magazine and Journal of Science
\textbf{42}(333), 1161--1177 (1951)

\bibitem{lakatos1976}
Lakatos, I.: Proofs and Refutations: The Logic of Mathematical Discovery.
Cambridge University Press (1976)

\bibitem{fagin1988belief}
Fagin, R., Halpern, J.Y.: Belief, awareness, and limited reasoning.
Artificial Intelligence \textbf{34}(1), 39--76 (1988)

\bibitem{wilson-private}
Wilson, R.A.: Private communication (2026)

\end{thebibliography}

\end{document}
